\uuid{HXRf}
\chapitre{Statistique}
\niveau{L2}
\module{Analyse de données}
\sousChapitre{Autre}
\titre{Simulation et classement sur tableur}
\theme{statistiques, tableur}
\auteur{}
\datecreate{2022-09-30}
\organisation{AMSCC}
\difficulte{}
\contenu{
\question{Simuler 200 tirages indépendants avec remise d'un élève parmi 13, chaque élève ayant la même probabilité. Représenter la répartition des résultats.}
\indication{Générer les numéros 1 à 13 puis compter le nombre d'apparitions de chacun.}
\reponse{\href{https://stcyrterrenetdefensegouvf-my.sharepoint.com/:x:/g/personal/maxime_nguyen_st-cyr_terre-net_defense_gouv_fr/EZAcxd2Zm4pIlrVf0mC7TyABy5-LvDvznM_4P6Amfs1m9A?e=480NAF}{Lien vers le fichier tableur}
\par
En A2 : \texttt{=ALEA.ENTRE.BORNES(1;13)}, puis recopier jusqu'à A201 pour obtenir 200 tirages. La fonction génère de nouveaux résultats lors d'un recalcul.
\par
Pour conserver une simulation, copier les résultats puis effectuer un collage spécial « valeurs ».
\par
Placer les numéros 1 à 13 en C2:C14. En D2 : \texttt{=NB.SI(\$A\$2:\$A\$201;C2)}, puis recopier jusqu'à D14 pour compter les apparitions de chaque élève.
\par
Représenter les effectifs par un diagramme en bâtons. Leur somme doit être 200. Chaque effectif a pour espérance
\[
\frac{200}{13}\simeq15{,}38,
\]
et les résultats fluctuent d'une simulation à l'autre.}
}
